% !TeX root = ../../Book2.tex
% 纲要标题：### 20.4 PBW 型思想
% 纲要位置：计划纲要.md，第 1409 行。

\section{PBW 型思想}
\label{b2:sec:2004}

Weyl 代数的最高次关系是交换关系，但它本身不交换。这提示我们先分析伴随分次，再逐次恢复较低次项。不过，存在从交换多项式环到伴随分次的满射，不等于这个满射一定单射；额外关系可能使预期的正规单项式彼此相关。

% 纲要要点（供目录核对）：
%
% * 由交换的伴随分次控制非交换代数
% * 与第五卷包络代数的 PBW 定理的联系
% * 形变不等于任意改动结构常数
%

\subsection{伴随分次与非交换结构}
\label{b2:subsec:2004:01}

\begin{proposition}\label{b2:prop:pbw-filtered-criterion}
设 \(A\) 由 \(x_1,\ldots,x_n\) 作为 \(k\) 代数生成，取字长滤过，并假设 \([x_i,x_j]\in F_1A\)。则有满同态
\[
k[X_1,\ldots,X_n]\longrightarrow\operatorname{gr}_F A,
\qquad X_i\longmapsto x_i+F_0A\in F_1A/F_0A.
\]
它是同构，当且仅当全部有序单项式 \(x_1^{a_1}\cdots x_n^{a_n}\) 是 \(A\) 的一组基。
\end{proposition}
\begin{proof}
交换子的次数小于二，所以一次符号交换；字长滤过又使这些符号生成整个伴随分次，得到满同态。由关系 \(x_jx_i=x_ix_j+[x_j,x_i]\)，按字长及错序数归纳，可以把任意长度不超过 \(m\) 的字整理为全次数不超过 \(m\) 的有序单项式的线性组合。

若这些单项式已经是一组基，就连各个滤过项的基也被上述整理确定。相应次数的商片以相同次数的单项式符号为基，故满同态逐次单射。

反过来，若伴随分次同构于多项式环，任意非零的有限有序单项式关系中，取最高全次数的部分，其符号就是交换多项式单项式之间的非平凡关系，矛盾。因此有序单项式线性无关；生成性已证，得到基。
\end{proof}

满足这种有序基性质时，常称所讨论的构造具有 PBW 型性质。这里所指的是一个已经说明的基与滤过结论，而不是把任意非交换关系都统称为 PBW 定理。

\ProseParagraphBreak
Weyl 代数给出完全可算的形变。它的 Rees 代数由 \(X=xz,D=dz,Z=z\) 生成，满足
\[
ZX=XZ,\quad ZD=DZ,\quad DX-XD=Z^2.
\]
这些关系给出整个代数。先将 \(Z\) 移到右侧，再用 \(DX=XD+Z^2\) 移动错序的 \(D,X\)；修正项使 \(X,D\) 的总个数减少二，因此按这个数和错序数归纳，可以把任何字写成 \(X^iD^jZ^r\) 的线性组合。它们在 Rees 代数中的像为 \(x^id^jz^{i+j+r}\)，由\cref{b2:thm:weyl-normal-basis}及不同 \(z\) 次数的独立性，它们线性无关。因此该关系代数与 Rees 代数同构。

\ProseParagraphBreak
它作为 \(k[Z]\) 模以 \(X^iD^j\) 为基。在 \(Z=0\) 处得到交换平面，在 \(Z=1\) 处得到 Weyl 代数；这是\cref{b2:thm:rees-fibres}的一个具体实现。参数 \(Z\) 及其平方不能随意从关系中删改，否则所得代数族未必保留这些基。

\subsection{包络代数与 PBW 定理}
\label{b2:subsec:2004:02}

为说明名称的来源，先给出所需对象。一个\term[Lie-daishu]{Lie 代数}[Lie algebra]是向量空间 \(\mathfrak g\)，带有双线性括号，满足 \([x,x]=0\) 及 Jacobi 恒等式
\[
\Bigl[x,[y,z]\Bigr]+\Bigl[y,[z,x]\Bigr]+\Bigl[z,[x,y]\Bigr]=0.
\]
其\term[fan-baoluo-daishu]{泛包络代数}[universal enveloping algebra]定义为
\[
U(\mathfrak g)=T_k(\mathfrak g)/\Bigl(xy-yx-[x,y]\mid x,y\in\mathfrak g\Bigr),
\]
分母表示这些元素生成的双边理想。它与\cref{b2:def:enveloping-algebra}的 \(A^e=A\otimes A^{\mathrm{op}}\) 是不同构造。

\ProseParagraphBreak
对有限维 \(\mathfrak g\) 选有序基 \(x_1,\ldots,x_n\)。关系的右侧为一次元素，所以字长滤过给出 \(\operatorname{Sym}(\mathfrak g)\to\operatorname{gr}U(\mathfrak g)\) 的自然满射。它的单射性是下面具名结果的内容。

\begin{claim}[PBW，证明见第五卷第15章]\label{b2:claim:pbw-enveloping-preview}
上述自然映射是同构；等价地，\(U(\mathfrak g)\) 以全部 \(x_1^{a_1}\cdots x_n^{a_n}\)、\(a_i\ge0\)，为一组基。
\end{claim}

这称为\term[PBW-dingli]{Poincaré--Birkhoff--Witt 定理}[Poincaré--Birkhoff--Witt theorem]。本节只陈述这项一般结果，完整证明归于第五卷第15章；当前章的 Weyl 基、Ore 正规式及后面的重写引理都给出独立证明，不把它当作前置。本节要强调的是：较低次的括号关系必须满足相容条件，Jacobi 恒等式在其中有实质作用。

\subsection{形变与结构常数的约束}
\label{b2:subsec:2004:03}

考虑三个生成元 \(x,y,z\)，随意指定
\[
yx=xy+x,\qquad zy=yz+y,\qquad zx=xz.
\]
这些关系的最高次部分都是交换关系，但不保证基仍是 \(x^ay^bz^c\)。

\begin{example}\label{b2:exm:incompatible-pbw-relations}
在上述关系定义的代数中，必有 \(x=0\)。因此它不具有预期的三变量 PBW 基。
\end{example}
\begin{proof}
对同一个字 \(zyx\) 先整理左边的 \(zy\)，得到
\[
(zy)x=(yz+y)x=yxz+yx=xyz+xz+xy+x.
\]
先整理右边的 \(yx\)，则得到
\[
z(yx)=z(xy+x)=xzy+xz=xyz+xy+xz.
\]
结合律使两者相等，故 \(x=0\)。这里出现的是一次额外关系，单看最高次的三条交换关系无法预见。
\end{proof}

这种失败可以用括号解释：指定的规则给出 \([y,x]=x\)、\([z,y]=y\)、\([z,x]=0\)，它们在以 \(x,y,z\) 为基的三维空间上不满足 Jacobi 恒等式。结合代数中的加法交换子却总满足该恒等式，展开三项后各个三重乘积成对消去；所以关系只能通过使 \(x\) 消失来相容。\secref{b2:sec:2006}将把这种同一字的两种整理方式作为重叠歧义处理。
